\section{问题重述}

帕金森病是一类常见神经退行性疾病，临床表现既包括静止性震颤、运动迟缓、肌强直等运动症状，也常伴随睡眠、认知、情绪和自主神经等非运动症状\upcite{clinical_pd}。除上述临床表现外，患者还可能出现发声稳定性下降、音强控制异常、音调波动和语音清晰度下降等语音障碍。已有研究表明，发声障碍和非侵入语音测试可用于帕金森病监测和辅助评估\upcite{little2009,tsanas2010}。语音采集成本较低，适合用于辅助筛查和持续随访，因此本题要求基于给定声学数据建立帕金森病特征识别与声音诊断模型。

题目包含五个相互关联的问题。第一问要求利用 Dataset 1 建立健康人与帕金森病患者的二分类辅助诊断模型；第二问要求识别与帕金森病判别相关的关键语音特征；第三问要求利用 Dataset 2 建立声音特征智能诊断模型；第四问希望区分运动型与非运动型帕金森病；第五问希望进一步区分六类临床症状。由于附件数据没有提供第四问所需的真实运动型/非运动型标签，也没有第五问所需的六类临床症状标签，本文不构建真实临床症状监督诊断模型，而是在 PD 受试者内部进行无监督探索，得到基于语音特征的二类和六类潜在表型。

本文工作的边界是基于既有数据和已确定的计算结果形成完整论文表达，不引入新的实验、不改变已有模型输出。所有模型结论均限于语音声学数据的统计关联和预测表现，不能解释为临床因果结论或临床确诊结论。
